\begin{lemma}
\label{lemma-relative-dualizing-flat-vanishing}
Let $R \to A$ be a flat finite type homomorphism of Noetherian rings.
Let $\mathfrak q \subset A$ be a prime ideal lying over
$\mathfrak p \subset R$. Then
$$
H^i(\omega_{A/R}^\bullet)_\mathfrak q \not = 0
\Rightarrow - d \leq i \leq 0
$$
where $d$ is the dimension of the fibre of $\Spec(A) \to \Spec(R)$
over $\mathfrak p$ at the point $\mathfrak q$. If all fibres of
$\Spec(A) \to \Spec(R)$ have dimension $\leq d$, then
$\omega_{A/R}^\bullet$ has tor amplitude in $[-d, 0]$
as a complex of $R$-modules.
\end{lemma}

\begin{proof}
The lower bound has been shown in
Lemma \ref{lemma-relative-dualizing-trivial-vanishing}.
Choose a factorization $R \to P \to A$ with $P = R[x_1, \ldots, x_n]$
as in Section \ref{section-relative-dualizing-complex-algebraic}
so that $\omega_{A/R}^\bullet = R\Hom(A, P)[n]$.
The upper bound means that $\Ext^i_P(A, P)$ is zero for $i > n$.
This follows from
More on Algebra, Lemma \ref{more-algebra-lemma-perfect-over-polynomial-ring}
which shows that $A$ is a perfect $P$-module with
tor amplitude in $[-n, 0]$.

\medskip\noindent
Proof of the final statement. Let $R \to R'$ be a ring homomorphism
of Noetherian rings. Set $A' = A \otimes_R R'$. Then
$$
\omega_{A'/R'}^\bullet =
\omega_{A/R}^\bullet \otimes_A^\mathbf{L} A' =
\omega_{A/R}^\bullet \otimes_R^\mathbf{L} R'
$$
The first isomorphism by Lemma \ref{lemma-base-change-relative-algebraic}
and the second, which takes place in $D(R')$, by
More on Algebra, Lemma \ref{more-algebra-lemma-base-change-comparison}.
By the first part of the proof
(note that the fibres of $\Spec(A') \to \Spec(R')$ have dimension $\leq d$)
we conclude that $\omega_{A/R}^\bullet \otimes_R^\mathbf{L} R'$
has cohomology only in degrees $[-d, 0]$. Taking $R' = R \oplus M$
to be the square zero thickening of $R$ by a finite $R$-module $M$,
we see that $R\Hom(A, P) \otimes_R^\mathbf{L} M$
has cohomology only in the interval $[-d, 0]$ for any finite $R$-module $M$.
Since any $R$-module is a filtered colimit of finite $R$-modules
and since tensor products commute with colimits we conclude.
\end{proof}